\subsection{例}

例 1. 设 \(\mathfrak{g}\) 是 \(\mathbf{K}\) 上的李代数，\(\mathbf{M}\) 是 \(\mathfrak{g}\)-模。\(\mathfrak{g}\) 上的 \(\mathbf{M}\)-模结构与 \(\mathfrak{g}\) 上的平凡 \(\mathbf{K}\)-模结构，在由 \(\mathfrak{g}\) 上的双线性形式组成的 \(\mathbf{K}\)-模 \(N = \mathcal{L}(M, M; K)\) 上定义了 \(\mathbf{M}\)-模结构。于是

\[
(x _ {N}. \beta) (m, m ^ {\prime}) = - \beta (x _ {M}. m, m ^ {\prime}) + \beta (m, x _ {M}. m ^ {\prime})\tag{10}
\]

对所有 \(x \in \mathfrak{g}\)、\(m, m'\) 属于 \(\mathbf{M}\)、\(\beta \in \mathbf{N}\) 成立。若 \(\beta\) 是 \(\mathbf{N}\) 的给定元素，则满足 \(x \in \mathfrak{g}\) 的 \(x_{\mathbf{N}}. \beta = 0\) 的集合是 \(\mathfrak{g}\) 的子代数。

设 \(\mathbf{M}\) 是 K-模，\(\beta\) 是 \(\mathbf{M}\) 上的双线性形式。由上面所述，满足 \(x \in \mathfrak{gl}(\mathbf{M})\) 的

\[
\beta (x. m, m ^ {\prime}) + \beta (m, x. m ^ {\prime}) = 0
\]

对所有 \(m \in \mathbf{M}\) 和 \(m' \in \mathbf{M}\)，这是 \(\mathfrak{gl}(\mathbf{M})\) 的李子代数。假设 \(\mathbf{K}\) 是域，\(\mathbf{M}\) 有限维且 \(\beta\) 非退化。则每个 \(x \in \mathfrak{gl}(\mathbf{M})\) 都有一个左伴随 \(x^*\)（相对于 \(\beta\)），它处处有定义；所述子代数正是满足 \(x \in \mathfrak{gl}(\mathbf{M})\) 的 \(x^* = -x\) 的集合。由此可以构造李代数的两个重要例子：

\begin{enumerate}
\def\labelenumi{(\alph{enumi})}
\tightlist
\item
  取 \(\mathbf{M} = \mathbf{K}^n\)，并
\end{enumerate}

\[
\beta ((\xi_ {1}, \dots , \xi_ {n}), (\eta_ {1}, \dots , \eta_ {n})) = \xi_ {1} \eta_ {1} + \dots + \xi_ {n} \eta_ {n}.
\]

我们把 \(\mathfrak{gl}(\mathbf{K}^n)\) 典范地等同于 \(\mathbf{M}_n(\mathbf{K})\)。于是得到的李代数就是斜对称矩阵的李代数。\(*\)（当 \(\mathbf{K} = \mathbf{R}\) 时，这个代数是正交群 \(\mathbf{O}(n, \mathbf{R}))\) 的李代数。）

\begin{enumerate}
\def\labelenumi{(\alph{enumi})}
\setcounter{enumi}{1}
\tightlist
\item
  取 \(\mathbf{M} = \mathbf{K}^{2m}\)，并
\end{enumerate}

\[
\beta ((\xi_ {1}, \dots , \xi_ {2 m}), (\eta_ {1}, \dots , \eta_ {2 m})) = \xi_ {1} \eta_ {m + 1} - \eta_ {1} \xi_ {m + 1} + \dots + \xi_ {m} \eta_ {2 m} - \eta_ {m} \xi_ {2 m}.
\]

相对于 \(\beta\) 的典范基，\(\mathbf{K}^{2m}\) 的矩阵是矩阵 \(\left( \begin{array}{cc}0 & I_m\\ -I_m & 0 \end{array} \right)\)。设 \(U = \left( \begin{array}{cc}A & B\\ C & D \end{array} \right)\) 是 \(\mathrm{K}^{2m}\) 中元素 \(u\) 相对于 \(\mathfrak{gl}(\mathbf{M})\) 的典范基的矩阵（\(A,B,C,D\) 属于 \(\mathbf{M}_m(\mathbf{K})\)）。根据《代数》第 IX 章 § 1 第 10 号的公式 (50)，\(u^*\) 相对于同一基的矩阵为

\[
\left( \begin{array}{c c} 0 & - I _ {m} \\ I _ {m} & 0 \end{array} \right) \left( \begin{array}{c c} ^ {t} A & ^ {t} C \\ ^ {t} B & ^ {t} D \end{array} \right) \left( \begin{array}{c c} 0 & I _ {m} \\ - I _ {m} & 0 \end{array} \right) = \left( \begin{array}{c c} ^ {t} D & - ^ {t} B \\ - ^ {t} C & ^ {t} A \end{array} \right).
\]

因此，条件 \(u^{*} = -u\) 等价于条件

\[
D = - ^ {t} A \qquad B = ^ {t} B \qquad C = ^ {t} C.
\]

*当 \(\mathbf{K} = \mathbf{R}\) 时，得到的李代数是辛群 \(\mathbf{Sp}(2m, \mathbf{R})\) 的李代数。*

例 2. 保持例 1 的记号。

M 上的 g-模结构在由 M 的自同态组成的 K-模 \(P = \mathcal{L}_{K}(M, M)\) 上定义了 g-模结构。由 (6)，对所有 \(x \in g\) 和 \(u \in P\)：

\[
x _ {\mathrm{P}}. u = \left[ x _ {\mathrm{M}}, u \right] = (\operatorname{ad} x _ {\mathrm{M}}). u\tag{11}
\]

其中 ad \(x_{M}\) 表示 \(x_{M}\) 在 \(\mathfrak{gl}(M)\) 的伴随表示下的像。换言之：

\[
x _ {\mathrm{P}} = \operatorname{ad} x _ {\mathrm{M}}\tag{12}
\]

在 \(\mathcal{L}(\mathcal{L}(\mathrm{M},\mathrm{M})) = \mathcal{L}(\mathfrak{gl}(\mathrm{M}))\) 中。

